\section{Conclusion}
\label{sec: conclusion}

Our work reconciles the long-standing trade-off between sublinear approximation and practical running time for GSTP through: (i)~\BasicPTS, which, via a holistic analysis of its two-step greedy approximation, achieves an $O(\sqrt{g})$~ratio; and (ii)~\ImprovPTS, which further incorporates a suspendable search mechanism and a pruning strategy that exploits the monotonicity and unimodality of the greedy function to match the efficiency of state-of-the-art linear-approximation solvers. This combined effectiveness and efficiency is validated by an extensive evaluation across five real-world datasets spanning tasks such as property search in road networks, team formation in social networks, and keyword-based exploration of KGs, where \ImprovPTS produces the highest-quality GSTs among approximation algorithms without compromising on speed.

While these results are promising, several important limitations point to clear avenues for advancement. First, regarding efficiency, our method is on par with the fastest existing index-free baselines but has not advanced beyond; its reliance on $g$~full SSSP computations remains a notable limitation at scale. Second, concerning approximation quality, our work is currently limited to the $d=2$ case of the $d$-star framework, despite the fact that larger~$d$ generally yields higher quality. Future directions, therefore, involve: (i)~developing lightweight, updatable indices (rather than the heavy, static one in KeyKG$^+$) to accelerate distance queries; and (ii)~deeply exploring the $d$-star structure to efficiently expand the search space.
% (ii)~tightening the theoretical analysis to reduce computational costs without compromising solution quality;

% Despite these encouraging results, there remains room for improvement. In particular, our practical running time does not yet surpass the speed bottleneck of the fastest index-free algorithms, and our algorithm still performs \(g\) full shortest-path computations, which may become costly when the number of groups is large~\cite{KGS_1}. A promising direction is to accelerate distance evaluation by introducing lightweight, fast-to-build, and dynamically maintainable indices, or by developing a tighter theoretical  characterization that further reduces the number of shortest-path computations without sacrificing solution quality.